% Part: lambda-calculus
% Chapter: syntax
% Section: substitution

\documentclass[../../../include/open-logic-section]{subfiles}

\begin{document}

\olfileid{lam}{syn}{sub}
\olsection{代入}

\begin{explain}
自由变元是对环境变元的引用，因此用特定的值去替换自由变元是有意义的。
例如，我们可能希望将 $\lambd[x][f x]$ 中的 $f$ 替换为一个具体的项，比如恒等函数 $\lambd[y][y]$。
结果为 $\lambd[x][(\lambd[y][y]) x]$。
将自由变元替换为 λ-项的过程称为\emph{代入}。
\end{explain}

\begin{defn}[代入]\ollabel{defn:substitution}
  将项 $N$ 代入项 $M$ 中的变元 $x$，记作 $\Subst{M}{N}{x}$，归纳定义如下：
  \begin{enumerate}
    \item $\Subst{x}{N}{x} = N$。\ollabel{defn:substitution-1}
    \item $\Subst{y}{N}{x}  = y$，若 $x \neq y$。\ollabel{defn:substitution-2}
    \item $\Subst{PQ}{N}{x}  = (\Subst{P}{N}{x}) (\Subst{Q}{N}{x})$。\ollabel{def:substitution-3}
    \item $\Subst{(\lambd[y][P])}{N}{x} = \lambd[y][\Subst{P}{N}{x}]$,
          若 $x \neq y$ 且 $y \notin \FV{N}$，否则无定义。\ollabel{defn:substitution-4}
  \end{enumerate}
\end{defn}

\begin{explain}
在 \olref{defn:substitution}\olref{defn:substitution-4} 中，我们要求 $x \neq y$，因为我们不希望用 $N$ 去替换 $M$ 中变元 $x$ 的\emph{约束}出现。
例如，如果我们要进行代入计算 $\Subst{\lambd[x][x]}{y}{x}$，结果不应该是 $\lambd[x][y]$，而应当是 $\lambd[x][x]$。

当在 $\lambd[y][P]$ 中将 $x$ 替换为 $N$ 时，我们还要求 $y \notin \FV{N}$。
例如，我们不能在 $\lambd[y][x]$ 中将 $x$ 替换为 $y$，即 $\Subst{\lambd[y][x]}{y}{x}$，因为那将得到 $\lambd[y][y]$，这个项表示的是接受一个自变元并直接将其返回的函数。
但项 $\lambd[y][x]$ 表示的是一个总是返回项 $x$（或 $x$ 所指代的内容）的函数。
所以，我们真正想要的结果是一个接受一个自变元、丢弃它、并返回环境变元 $y$ 的函数。
要正确地做到这一点，我们必须先“重命名”约束变元 $y$。
\end{explain}

\begin{prob}
  下列代入的结果是什么？
  \begin{enumerate}
  \item $\Subst{\lambd[y][x(\lambd[w][vwx])]}{(uv)}{x}$
  \item $\Subst{\lambd[y][x(\lambd[x][x])]}{(\lambd[y][xy])}{x}$
  \item $\Subst{y(\lambd[v][xv])}{(\lambd[y][vy])}{x}$
  \end{enumerate}
\end{prob}

\begin{thm}\ollabel{thm:notinfv}
  若 $x \notin \FV{M}$，则 $\FV{\Subst{M}{N}{x}} = \FV{M}$，前提是等式左端有定义。
\end{thm}

\begin{proof}
  对 $M$ 的构造进行归纳。
  \begin{enumerate}
  \item $M$ 是变元：练习。
  \item $M$ 形如 $(PQ)$：练习。
  \item $M$ 形如 $\lambd[y][P]$，且由于 $\Subst{\lambd[y][P]}{N}{x}$ 有定义，它必然是 $\lambd[y][\Subst{P}{N}{x}]$。则 $\Subst{P}{N}{x}$ 亦必有定义；同时，$x \neq y$ 且 $x \notin \FV{Q}$。于是：
    \begin{multline*}
      \FV{\Subst{\lambd[y][P]}{N}{x}} = \\
    \begin{aligned}
      & = \FV{\lambd[y][\Subst{P}{N}{x}]} &&
      \text{由 \olref{defn:substitution-4}}\\
      & = \FV{\Subst{P}{N}{x}} \setminus \{y\} &&
      \text{由 \olref[fv]{def:fv}\olref[fv]{def:fv2}}\\
      & = \FV{P} \setminus \{y\} && \text{由归纳假设} \\
      & = \FV{\lambd[y][P]} && \text{由 \olref[fv]{def:fv}\olref[fv]{def:fv2}}
    \end{aligned}
    \end{multline*}
  \end{enumerate}
\end{proof}

\begin{prob}
  完成 \olref[lam][syn][sub]{thm:notinfv} 的证明。
\end{prob}

\begin{thm}\ollabel{thm:infv}
  若 $x \in \FV{M}$，则 $\FV{\Subst{M}{N}{x}} = (\FV{M} \setminus
  \{x\}) \cup \FV{N}$，前提是等式左端有定义。
\end{thm}

\begin{proof}
  对 $M$ 的构造进行归纳。
  \begin{enumerate}
  \item $M$ 是变元：练习。
  \item $M$ 形如 $PQ$：由于 $\Subst{(PQ)}{N}{y}$ 有定义，它必然是 $(\Subst{P}{N}{x})(\Subst{Q}{N}{x})$，且这两个代入均有定义。此外，由于 $x \in \FV{PQ}$，可知 $x \in \FV{P}$ 或 $x \in \FV{Q}$ 或两者皆有。余下部分留作练习。
  \item $M$ 形如 $\lambd[y][P]$。由于 $\Subst{\lambd[y][P]}{N}{x}$ 有定义，它必然是 $\lambd[y][\Subst{P}{N}{x}]$，其中 $\Subst{P}{N}{x}$ 有定义，$x \neq y$ 且 $y \notin \FV{N}$；同时，由于 $y \in
    \FV{\lambd[x][P]}$，可知 $y \in \FV{P}$。现在：
    \begin{multline*}
      \FV{\Subst{(\lambd[y][P])}{N}{x}} = \\
      \begin{aligned}
      & = \FV{\lambd[y][\Subst{P}{N}{x}]} \\
      & = \FV{\Subst{P}{N}{x}} \setminus \{y\} \\
      & = ((\FV{P} \setminus \{y\}) \cup (\FV{N} \setminus \{x\})
       && \text{由归纳假设} \\
      & = (\FV{P} \setminus \{x, y\}) \cup \FV{N}
       && x \notin \FV{N} \\
      & = (\FV{\lambd[y][P]} \setminus \{x\}) \cup \FV{N}
      \end{aligned}
      \end{multline*}
  \end{enumerate}
\end{proof}

\begin{prob}
  完成 \olref[lam][syn][sub]{thm:infv} 的证明。
\end{prob}

\begin{thm}\ollabel{thm:clr}
  $x \notin \FV{\Subst{M}{N}{x}}$，前提是等式右端有定义且 $x \notin \FV{N}$。
\end{thm}

\begin{proof}
  练习。
\end{proof}

\begin{prob}
  证明 \olref[lam][syn][sub]{thm:clr}。
\end{prob}


\begin{thm}\ollabel{thm:inv}
  若 $\Subst{M}{y}{x}$ 有定义且 $y \notin \FV{M}$，则 $\Subst{\Subst{M}{y}{x}}{x}{y} = M$。
\end{thm}

\begin{proof}
  对 $M$ 的构造进行归纳。
  \begin{enumerate}
  \item $M$ 是变元 $z$：练习。
  \item $M$ 形如 $(PQ)$。则：
    \begin{align*}
      \Subst{\Subst{(PQ)}{y}{x}}{x}{y}
      &=\Subst{((\Subst{P}{y}{x})(\Subst{Q}{y}{x}))}{x}{y} \\
      &= (\Subst{\Subst{P}{y}{x}}{x}{y})(\Subst{\Subst{Q}{y}{x}}{x}{y}) \\
      &= (PQ) \text{ 由归纳假设}
    \end{align*}
  \item $M$ 形如 $\lambd[z][N]$。因为 $\Subst{\lambd[z][N]}{y}{x}$ 有定义，可知 $z \neq y$。所以：
    \begin{align*}
      \Subst{\Subst{(\lambd[z][N])}{y}{x}}{x}{y}\\
      & = \Subst{(\lambd[z][\Subst{N}{y}{x}])}{x}{y} \\
      & = \lambd[z][\Subst{\Subst{N}{y}{x}}{x}{y}] \\
      &= \lambd[z][N] \text{ 由归纳假设} 
    \end{align*}
  \end{enumerate}
\end{proof}

\begin{prob}
  完成 \olref[lam][syn][sub]{thm:inv} 的证明。
\end{prob}

\end{document}